\documentclass[11pt]{amsart}
\input{glyphtounicode}
\pdfgentounicode=1

\usepackage[T1]{fontenc}
\usepackage[utf8]{inputenc}
\usepackage{amsmath,amssymb,amsthm,mathtools}
\usepackage{array,booktabs,tabularx}
\usepackage{graphicx}
\usepackage{xcolor}
\usepackage{geometry}
\usepackage{enumitem}
\usepackage{hyperref}

\geometry{margin=1in}
\setlength{\emergencystretch}{3em}
\allowdisplaybreaks
\hypersetup{
  colorlinks=true,
  linkcolor=blue!45!black,
  citecolor=blue!45!black,
  urlcolor=blue!45!black,
  pdftitle={Kernel-Integrated Typed Multilinear Laws: A Finite Core, a Conditioned Analytic Extension, and an Open Geometric Program},
  pdfauthor={Eliuth Chavero Jasso},
  pdfsubject={Typed multilinear laws, recursive orthogonal projection, nodewise error certificates, kernel realizations, truncated Hodge/spectral extensions},
  pdfkeywords={multilinear tree, recursive projection, closure residual, nodewise certificate, kernel operator, Hodge theory, spectral dimension}
}

\newtheorem{theorem}{Theorem}[section]
\newtheorem{proposition}[theorem]{Proposition}
\newtheorem{corollary}[theorem]{Corollary}
\newtheorem{lemma}[theorem]{Lemma}
\theoremstyle{definition}
\newtheorem{definition}[theorem]{Definition}
\newtheorem{example}[theorem]{Example}
\theoremstyle{remark}
\newtheorem{remark}[theorem]{Remark}
\newtheorem{openproblem}[theorem]{Open Problem}

\DeclareMathOperator{\ran}{ran}
\DeclareMathOperator{\Int}{Int}
\DeclareMathOperator{\Tr}{Tr}
\DeclareMathOperator{\op}{op}

\title{Kernel-Integrated Typed Multilinear Laws:\\
A Finite Core, a Conditioned Analytic Extension,\\
and an Open Geometric Program}
\author{Eliuth Chavero Jasso}
\date{\today}

\begin{document}
\maketitle
\markboth{Eliuth Chavero Jasso}{Kernel-Integrated Typed Multilinear Laws}

\begin{abstract}
We give a self-contained formalization of SEION / Kernel-Integrated Laws
at three explicitly separated levels: a \emph{proved finite core}
(typed multilinear laws, recursive orthogonal projection, and exact
nodewise error propagation on finite composition trees), a
\emph{conditioned analytic extension} (kernel realizations on
$L^2$ spaces, truncated Hodge/cohomological compatibility, an induced
Markov Laplacian), and an \emph{open geometric program} (continuum
limits, microlocal regularity, algebraization) that is stated but not
claimed. The central proved result is the ambient-to-projected
coefficient improvement
$E_T^{\mathrm{amb}} \le k\rho M^{k-1}L_T$ versus
$E_T^{P} \le (k-1)\rho M^{k-1}L_T$
for a finite typed composition tree with $k$ internal nodes, proved by
an exact nodewise error-orthogonality identity together with induction.
We are explicit throughout about what is proved unconditionally, what
is proved only under stated hypotheses, what is demonstrated only in a
bounded or truncated regime, and what remains open. No claim in this
document is asserted stronger than the evidence in
\S\S\ref{sec:xxvi} supports.
\end{abstract}

\tableofcontents

\section{Scope and epistemic map}

The canonical formulation separates three levels:
\[
\underbrace{\text{finite proved core}}_{\S\S\ref{sec:types}--\ref{sec:xiii}}
\;\subset\;
\underbrace{\text{conditioned analytic extension}}_{\S\S\ref{sec:xv}--\ref{sec:xxiii}}
\;\subset\;
\underbrace{\text{open geometric program}}_{\S\ref{sec:xxiv}}.
\]
The scientifically most defensible object here is not a physical
theory, but a theory of multilinear laws, projection-based reduction,
and structured error propagation.

\section{Foundational algebraic data}
\label{sec:types}

\subsection{Type system}

Let $\mathcal T$ be a finite set of types. To each $\tau \in \mathcal T$
we associate a finite-dimensional Hilbert space $V_\tau$ over
$\mathbb K \in \{\mathbb R, \mathbb C\}$. A \emph{typed signature} is a
list $(\tau_1,\ldots,\tau_a;\tau_0)$ representing an operation with
inputs in $V_{\tau_1},\ldots,V_{\tau_a}$ and output in $V_{\tau_0}$.

\subsection{Typed $a$-ary law}

A \emph{typed multilinear law} is a map
$\mu : V_{\tau_1}\times\cdots\times V_{\tau_a} \to V_{\tau_0}$,
multilinear in each argument. In orthonormal bases it is determined by
a tensor
$K \in V_{\tau_0}\otimes V_{\tau_1}^*\otimes\cdots\otimes V_{\tau_a}^*$,
with coordinates
\[
\bigl[\mu(x_1,\ldots,x_a)\bigr]^{i_0}
= \sum_{i_1,\ldots,i_a} K^{i_0}_{\ i_1\cdots i_a}\prod_{j=1}^{a} x_j^{i_j}.
\]
Its operator norm is
$\lVert\mu\rVert_{\op} = \sup_{\lVert x_1\rVert=\cdots=\lVert x_a\rVert=1}\lVert \mu(x_1,\ldots,x_a)\rVert$.
The computational formalization accepts dense, sparse, complex, typed,
and CP-factorized laws, with arity and domains tracked explicitly.

\section{CP representation}

To avoid storing a tensor of size $d_0 d_1\cdots d_a$, a rank-$R$ CP
representation is used:
\[
K = \sum_{r=1}^{R}\lambda_r\, o_r\otimes a_r^{(1)}\otimes\cdots\otimes a_r^{(a)},
\qquad
\mu(x_1,\ldots,x_a) = O\Bigl[\lambda\odot A^{(1)}x_1\odot\cdots\odot A^{(a)}x_a\Bigr],
\]
with $A^{(j)}:V_{\tau_j}\to\mathbb K^R$ and $O:\mathbb K^R\to V_{\tau_0}$.

The CP representation is not unique: for nonzero scalars $c_r^{(j)}$
with $c_r^{(0)}c_r^{(1)}\cdots c_r^{(a)}=1$, factors may be rescaled
without changing $K$, and there is additional freedom under
permutation of components. Comparison of factors must therefore be
performed modulo scale $+$ phase $+$ permutation $+$ admissible gauge
transformations.

\begin{remark}
CP factorization is a computational representation; on its own it does
not constitute a new algebraic structure and does not by itself prove
identifiability.
\end{remark}

\section{Symmetries and identities}

\subsection{Cyclic symmetry}

For an $a$-ary law, cyclicity means
$\mu(x_1,\ldots,x_a) = \mu(x_2,\ldots,x_a,x_1)$.
The cyclic projector on the tensor space is
$\Pi_{\mathrm{cyc}} = \tfrac1a\sum_{j=0}^{a-1}\sigma^j$,
where $\sigma$ rotates the input slots. Since $\sigma$ acts unitarily,
$\Pi_{\mathrm{cyc}}^2=\Pi_{\mathrm{cyc}}$ and
$\Pi_{\mathrm{cyc}}^*=\Pi_{\mathrm{cyc}}$, so
$\mu_{\mathrm{cyc}} := \Pi_{\mathrm{cyc}}\mu$ is cyclic \emph{by
construction}.

\begin{remark}
A zero cyclic residual after this operation is a structural identity,
not evidence that the symmetry was learned.
\end{remark}

\subsection{Filippov's fundamental identity}

For a ternary law $\mu$, one convention for the Filippov residual is
\begin{align*}
\operatorname{FI}_\mu(x_1,x_2;y_1,y_2,y_3)
={}&\mu(x_1,x_2,\mu(y_1,y_2,y_3))\\
&-\mu(\mu(x_1,x_2,y_1),y_2,y_3)\\
&-\mu(y_1,\mu(x_1,x_2,y_2),y_3)\\
&-\mu(y_1,y_2,\mu(x_1,x_2,y_3)),
\end{align*}
and the Filippov condition is $\operatorname{FI}_\mu = 0$.

\begin{remark}
FI, GJI, Jacobi, cyclicity, and associativity are distinct identities.
A small residual for one does not imply the others.
\end{remark}

\section{Associator conventions}

For a ternary operation there is no unique canonical associator
without specifying how compositions are inserted.

\subsection{Five-input ternary associator}
\[
A_\mu^{(5)}(x_1,\ldots,x_5)
= \mu(\mu(x_1,x_2,x_3),x_4,x_5) - \mu(x_1,x_2,\mu(x_3,x_4,x_5)),
\qquad A_\mu^{(5)}: V^5\to V.
\]

\subsection{Anchored binary reduction}

Fixing an anchor $e\in V$, set $x\circ_e y := \mu(x,y,e)$. The induced
binary associator is
$A_e(x,y,z) = (x\circ_e y)\circ_e z - x\circ_e(y\circ_e z)$.

\subsection{Operadic insertions}

For an $a$-ary law, insertion at slot $i$ is $\mu\circ_i\mu$. In the
ternary case,
$(\mu\circ_1\mu)(x_1,\ldots,x_5)=\mu(\mu(x_1,x_2,x_3),x_4,x_5)$ and
$(\mu\circ_3\mu)(x_1,\ldots,x_5)=\mu(x_1,x_2,\mu(x_3,x_4,x_5))$, so
$A_\mu^{(5)} = \mu\circ_1\mu - \mu\circ_3\mu$.

\begin{remark}
These conventions have different domains and meanings and must not be
merged under a single ``associator score.''
\end{remark}

\section{Curvature: definition versus theorem}

\subsection{Constitutive curvature}

One may \emph{define}, by convention, $R_{\mathrm{alg}} := A_\mu$. This
is a definition of algebraic curvature, not an automatic identification
with the standard curvature of a connection.

\subsection{Curvature induced by left operators}

For a binary product $\circ$, let $L_x z = x\circ z$ and
$[x,y]=x\circ y - y\circ x$. Define
$R_{\mathrm{std}}(x,y) := [L_x,L_y]-L_{[x,y]}$.

\begin{theorem}
\label{thm:curvature}
With $A(x,y,z) = (x\circ y)\circ z - x\circ(y\circ z)$,
\[
R_{\mathrm{std}}(x,y)z = A(y,x,z) - A(x,y,z).
\]
\end{theorem}
\begin{proof}
\begin{align*}
R_{\mathrm{std}}(x,y)z
&= x\circ(y\circ z) - y\circ(x\circ z) - (x\circ y)\circ z + (y\circ x)\circ z\\
&= \bigl[x\circ(y\circ z) - (x\circ y)\circ z\bigr]
 - \bigl[y\circ(x\circ z) - (y\circ x)\circ z\bigr]\\
&= -A(x,y,z) + A(y,x,z). \qedhere
\end{align*}
\end{proof}

\begin{remark}
An outright identity $R_{\mathrm{std}}=A$ requires additional
hypotheses, particular symmetries, or a different convention. The
universal identity currently defensible is the antisymmetrized
difference of Theorem~\ref{thm:curvature}.
\end{remark}

\section{Projectors and the reduced law}

\subsection{Reduced spaces}

For each type $\tau$, let $W_\tau$ be another finite-dimensional
Hilbert space and $Q_\tau : W_\tau \to V_\tau$ an isometry,
$Q_\tau^*Q_\tau = I_{W_\tau}$. Define the orthogonal projector
$P_\tau = Q_\tau Q_\tau^*$, which satisfies
$P_\tau^2=P_\tau$, $P_\tau^*=P_\tau$, $\lVert P_\tau\rVert=1$.

\subsection{Reduced law}

For $\mu : V_{\tau_1}\times\cdots\times V_{\tau_a}\to V_{\tau_0}$, the
\emph{reduced law} is
\[
\bar\mu(z_1,\ldots,z_a) := Q_{\tau_0}^{*}\,\mu(Q_{\tau_1}z_1,\ldots,Q_{\tau_a}z_a),
\]
with ambient realization
$Q_{\tau_0}\bar\mu(z_1,\ldots,z_a) = P_{\tau_0}\,\mu(Q_{\tau_1}z_1,\ldots,Q_{\tau_a}z_a)$.

\subsection{Exact closure}

The subspace family is invariant under $\mu$ if
$\mu(\ran P_{\tau_1},\ldots,\ran P_{\tau_a}) \subseteq \ran P_{\tau_0}$,
equivalently
$(I-P_{\tau_0})\,\mu(P_{\tau_1}x_1,\ldots,P_{\tau_a}x_a)=0$ for all
$x_i$.

\subsection{Closure leakage}

The leakage map is
$\mathcal C_\mu := (I-P_{\tau_0})\,\mu(P_{\tau_1}\cdot,\ldots,P_{\tau_a}\cdot)$,
with norm $\rho_\mu := \lVert\mathcal C_\mu\rVert_{\op}$, so
\[
\bigl\lVert (I-P_{\tau_0})\,\mu(P_{\tau_1}x_1,\ldots,P_{\tau_a}x_a)\bigr\rVert
\le \rho_\mu \prod_{i=1}^{a}\lVert x_i\rVert.
\]
Experiments also use a normalized diagnostic
\[
\mathcal E_{\mathrm{closure}} := \mathbb E\left[\frac{\lVert(I-P)\mu(PX_1,\ldots,PX_a)\rVert^2}
{\lVert\mu(PX_1,\ldots,PX_a)\rVert^2+\varepsilon}\right].
\]

\begin{remark}
The operator norm $\rho_\mu$ and the sampled average
$\mathcal E_{\mathrm{closure}}$ are not the same quantity; the latter
does not automatically certify the former.
\end{remark}

\section{Typed composition trees}

This is currently the strongest mathematical core of the theory. Let
$T$ be a finite, ordered, rooted tree. Each internal node $v$ carries
an arity $a_v$, an output type $\tau(v)$, input types
$\tau(v,1),\ldots,\tau(v,a_v)$, and a law
$\mu_v : \prod_{j=1}^{a_v}V_{\tau(v,j)} \to V_{\tau(v)}$. Edges must
respect types; an ill-typed tree carries no mathematical semantics.
Write $k(T) = \lvert \Int T\rvert$ for the number of internal nodes.

\section{Two evaluations of the tree}

\subsection{Ambient evaluation}

For a leaf $\ell$ with reduced datum $z_\ell \in W_{\tau(\ell)}$, set
$F_\ell = Q_{\tau(\ell)}z_\ell$. At an internal node,
$F_v = \mu_v(F_{c_1},\ldots,F_{c_{a_v}})$: no intermediate result is
projected.

\subsection{Recursively projected evaluation}

At leaves, $R_\ell = Q_{\tau(\ell)}z_\ell$; at each internal node,
\[
R_v := P_{\tau(v)}\,\mu_v(R_{c_1},\ldots,R_{c_{a_v}}).
\]
Define the nodewise error $D_v = F_v - R_v$ and its components
$D_v^P = P_v D_v$, $D_v^N = (I-P_v)D_v$, separating retained error from
normal leakage.

\section{Exact root-error geometry}
\label{sec:root-geometry}

Let $r$ be the root; write $F=F_r$, $R=R_r$, $P=P_r$, and define
\[
E_T^{\mathrm{amb}} = \lVert F-R\rVert,\quad
E_T^{P} = \lVert PF-R\rVert,\quad
E_T^{N} = \lVert (I-P)F\rVert,\quad
E_T^{\mathrm{red}} = \lVert Q^*F - Q^*R\rVert.
\]

\begin{theorem}[Error orthogonality]
\label{thm:orthogonality}
$(E_T^{\mathrm{amb}})^2 = (E_T^P)^2 + (E_T^N)^2$ and
$E_T^{\mathrm{red}} = E_T^P$.
\end{theorem}
\begin{proof}
Since $R\in\ran P$, $F-R = (PF-R)+(I-P)F$, and the two terms are
orthogonal (one lies in $\ran P$, the other in its complement), giving
the first identity. Since $Q^*(I-P)=0$ and $Q^*$ is isometric on
$\ran P$, $Q^*F-Q^*R = Q^*(PF-R)$ has the same norm as $PF-R$, giving
the second.
\end{proof}

The reduced system observes exactly the projected component of the
ambient error.

\section{Exact subset expansion}

Let $v$ have arity $a$ and write $F_i=R_i+D_i$ for its children. By
multilinearity, for each nonempty $S\subseteq[a]$ define
$y_i^S = D_i$ if $i\in S$, else $R_i$, and let the local residual be
$r_v(R_1,\ldots,R_a) = (I-P_v)\mu_v(R_1,\ldots,R_a)$.

\begin{theorem}[Exact nodewise expansion]
\label{thm:subset}
\[
D_v = r_v(R_1,\ldots,R_a) + \sum_{\varnothing\neq S\subseteq[a]} \mu_v(y_1^S,\ldots,y_a^S).
\]
\end{theorem}
This identity is exact. Terms with $\lvert S\rvert=1$ are single-branch
propagation; $\lvert S\rvert\ge 2$ terms represent genuine
cross-branch interaction. Since $P_v r_v = 0$,
\[
D_v^P = \sum_{\varnothing\neq S\subseteq[a]} P_v\mu_v(y_1^S,\ldots,y_a^S).
\]
A node's own local leakage does not appear in its own projected error,
but a child's normal leakage can become tangential output,
$P_v\mu_v(\ldots,D_i^N,\ldots)\neq 0$ — the \emph{normal-to-tangent
mechanism}.

\section{Mixed-mask calculus}

Each input slot of a node may be in one of three states
$\mathfrak S=\{R,P,N\}$. For a mask
$\sigma=(\sigma_1,\ldots,\sigma_a)\in\mathfrak S^a$, set
$Z_i^\sigma = R_i$ if $\sigma_i=R$, $D_i^P$ if $\sigma_i=P$, $D_i^N$ if
$\sigma_i=N$. Then
\[
\mu_v(F_1,\ldots,F_a) = \sum_{\sigma\in\{R,P,N\}^a}\mu_v(Z_1^\sigma,\ldots,Z_a^\sigma).
\]
The mask $(R,\ldots,R)$ contains the reduced evaluation and the local
residual; other masks contain error propagation and interaction.

For $\chi\in\{P,N\}$, block gains
$\beta_v^\chi(\sigma) = \lVert \Pi_v^\chi\,\mu_v(\Pi_{c_1}^{\sigma_1}\cdot,\ldots,\Pi_{c_a}^{\sigma_a}\cdot)\rVert_{\op}$
(with $\Pi^R=\Pi^P=P$, $\Pi^N=I-P$, though $R$ and $P$ states retain
different magnitude bounds) satisfy the recurrence
\[
b_v^\chi \le \lambda_v^\chi + \sum_{\sigma\neq(R,\ldots,R)}\beta_v^\chi(\sigma)\prod_i b_i^{\sigma_i},
\qquad \lambda_v^P=0,\quad \lambda_v^N\le\rho_v\prod_i b_i^R.
\]
This recurrence is a finer nodewise certificate than a single global
norm; its soundness is proved for the declared finite configuration.

\section{Scalar telescoping and optimal order}

To compare $\mu(F_1,\ldots,F_a)$ against $\mu(R_1,\ldots,R_a)$, slots
are swapped one at a time. With $e_i$ an error bound, $r_i,f_i$ bounds
on $\lVert R_i\rVert,\lVert F_i\rVert$, $G_i$ a slot gain, and
$w_i=G_ie_i$, the cost of order $\pi$ is
\[
C(\pi) = \sum_{t=1}^{a} w_{\pi_t}\prod_{s<t}r_{\pi_s}\prod_{s>t}f_{\pi_s}.
\]

\begin{proposition}
Adjacent transposition favors placing $i$ before $j$ iff
$w_i(f_j-r_j) \le w_j(f_i-r_i)$.
\end{proposition}
Ordering by this rule minimizes $C(\pi)$ within this declared scalar
certificate family — not among all conceivable bounding methods.

\section{The structural $k\to k-1$ theorem}
\label{sec:xiii}

\subsection{Uniform hypotheses}

Suppose $\lVert\mu_v\rVert_{\op}\le M$ for every node, and
\[
\lVert (I-P_v)\mu_v(P_{c_1}\cdot,\ldots,P_{c_a}\cdot)\rVert_{\op}\le\rho.
\]
Let $L_T = \prod_{\ell\in\text{Leaves}(T)}\lVert z_\ell\rVert$.

\begin{theorem}[Ambient-to-projected coefficient improvement]
\label{thm:kk1}
For every finite, well-typed tree with $k$ internal nodes,
\[
E_T^{\mathrm{amb}} \le k\rho M^{k-1}L_T,
\qquad
E_T^{N} \le k\rho M^{k-1}L_T,
\qquad
E_T^{P} = E_T^{\mathrm{red}} \le (k-1)\rho M^{k-1}L_T.
\]
\end{theorem}
\begin{proof}
Let the root have arity $a$, subtrees $T_i$, $k_i=k(T_i)$, so
$k = 1+\sum_i k_i$. By induction, $\lVert F_i-R_i\rVert \le k_i\rho M^{k_i-1}L_i$
and $\lVert F_i\rVert,\lVert R_i\rVert \le M^{k_i}L_i$. Telescoping over
children,
\begin{align*}
\lVert\mu(F_1,\ldots,F_a)-\mu(R_1,\ldots,R_a)\rVert
&\le \sum_i M\cdot k_i\rho M^{k_i-1}L_i\prod_{j\neq i}M^{k_j}L_j\\
&= \Bigl(\sum_i k_i\Bigr)\rho M^{k-1}L_T = (k-1)\rho M^{k-1}L_T.
\end{align*}
The root's own new residual satisfies
$\lVert(I-P)\mu(R_1,\ldots,R_a)\rVert \le \rho M^{k-1}L_T$, so
$E_T^{\mathrm{amb}} \le (k-1)\rho M^{k-1}L_T + \rho M^{k-1}L_T = k\rho M^{k-1}L_T$.
But projecting the root gives $P(I-P)\mu(R_1,\ldots,R_a)=0$: the
root's own local residual vanishes under projection, leaving
$E_T^{P}\le (k-1)\rho M^{k-1}L_T$.
\end{proof}

\begin{remark}
The improvement does not depend on tree balance or any statistical
distribution; it comes from the orthogonality of the final projection
(Theorem~\ref{thm:orthogonality}).
\end{remark}

\subsection{Base cases}

For $k=1$: $F=\mu(Qz_1,\ldots,Qz_a)$, $R=P\mu(Qz_1,\ldots,Qz_a)$, so
$PF-R=0$ and $E_T^P=E_T^{\mathrm{red}}=0$ exactly, though
$E_T^{\mathrm{amb}}$ may be nonzero. For $k=2$,
$T(x,y,z)=\mu_v(\mu_u(x,y),z)$: $E_T^{\mathrm{amb}}\le 2\rho ML_T$ but
$E_T^{P}\le \rho ML_T$ — leakage created at $u$ can be turned tangential
by $v$; leakage created at $v$ itself vanishes on projection.

\section{Path-sum certificate}

Given a valid scalar recurrence $B_v \le \lambda_v + \sum_j h_{v,j}B_{c_j}$
(with $\lambda_v$ a local residual source and $h_{v,j}$ a transport
gain), recursive substitution gives
\[
B_{T,\mathrm{path}}^{\mathrm{amb}}
= \sum_{v\in\Int T} \lambda_v \prod_{(a,j)\in\mathrm{path}(v,\mathrm{root})} h_{a,j},
\qquad
B_{T,\mathrm{path}}^{P}
= \sum_{\substack{v\in\Int T\\ v\neq\mathrm{root}}} \lambda_v \prod_{(a,j)\in\mathrm{path}(v,\mathrm{root})} \tilde h_{a,j},
\]
where the root's source is omitted in the projected version and the
last factor uses the projected output gain. This attributes error to
each node and path individually rather than reporting only one global
constant.

\section{Signed forests and identities}
\label{sec:signed-forests}

For a finite compatible combination $\mathcal F = \sum_{\alpha=1}^m c_\alpha T_\alpha$,
the triangle inequality gives
$E_{\mathcal F}^{P} \le \sum_\alpha \lvert c_\alpha\rvert B_{T_\alpha}^P$.
For the five-input ternary associator $A_\mu^{(5)}=T_1-T_2$ with
$k(T_1)=k(T_2)=2$, each contributes at most $(k-1)\rho ML=\rho ML$, so
$\lvert A_\mu^{(5),P}\rvert \le 2\rho ML$ (the ambient triangle bound
would be $4\rho ML$).

\begin{openproblem}
The coefficient $2$ is an immediate bound; its sharpness, and whether
shared-DAG or adversarial whole-forest cancellation can lower it, are
open (see \S\ref{sec:xxvi}).
\end{openproblem}

\section{Representation error}

Suppose each $\mu_v$ is replaced by an approximation $\widehat\mu_v$
with $\lVert\mu_v\rVert\le M$, $\lVert\mu_v-\widehat\mu_v\rVert\le\delta$,
$\lVert\widehat\mu_v\rVert\le M+\delta$. For a homogeneous tree of $k$
nodes,
\[
E_{\mathrm{repr}} \le k\delta(M+\delta)^{k-1}L_T,
\]
separating law-approximation error from recursive-projection error: a
CP or quantized representation introduces a term distinct from closure
leakage.

\section{Extremal constants}

With $\eta=\rho/M$, the optimal projected constant for a tree $T$ is
\[
C_T^{P}(\eta) := \sup_{\mathcal A(T,\eta)} \frac{E_T^P}{\rho M^{k-1}L_T},
\]
over the admissible class $\mathcal A(T,\eta)$. Theorem~\ref{thm:kk1}
proves $C_T^P(\eta)\le k-1$ but not $C_T^P(\eta)=k-1$. Open
possibilities: equality, strict inequality, or a finer dependence
$C_T^P(\eta) = C(T,\eta,\dim V,\operatorname{rank}P,\text{arities})$.
Current evidence includes lower constructions and certified small
cells, but sharpness at $\eta>0$ remains open.

\section{Kernel-integrated realization}
\label{sec:xv}

\subsection{Function-space setting}

Let $(X,\mu)$ be a $\sigma$-finite measure space and $H=L^2(X,\mu)$.

\subsection{Ternary kernel}

For $K:X^4\to\mathbb C$, define
\[
T_K(f,g,h)(p) = \iiint_{X^3} K(p;q,r,s)f(q)g(r)h(s)\,d\mu(q)\,d\mu(r)\,d\mu(s).
\]

\begin{proposition}
If $K\in L^2(X^4)$, then $T_K:H^3\to H$ is bounded with
$\lVert T_K(f,g,h)\rVert_2 \le \lVert K\rVert_{L^2(X^4)}\lVert f\rVert_2\lVert g\rVert_2\lVert h\rVert_2$.
\end{proposition}
\begin{proof}
Cauchy--Schwarz on $X^3$ followed by Fubini in the output variable.
\end{proof}
The $n$-ary construction uses $K_n\in L^2(X^{n+1})$ and gives
$\lVert T_{K_n}(f_1,\ldots,f_n)\rVert_2 \le \lVert K_n\rVert_2\prod_j\lVert f_j\rVert_2$.

\subsection{Stability in the kernel}

If $K,\widetilde K\in L^2(X^4)$,
$\lVert T_K(f,g,h)-T_{\widetilde K}(f,g,h)\rVert_2 \le \lVert K-\widetilde K\rVert_2\lVert f\rVert_2\lVert g\rVert_2\lVert h\rVert_2$,
so $K_n\to K$ in $L^2$ implies strong convergence on each fixed triple.

\section{Kernel of the associator}

Two composite kernels arise for five inputs:
\[
K_L(p;q_1,\ldots,q_5) = \int_X K(p;u,q_4,q_5)K(u;q_1,q_2,q_3)\,d\mu(u),
\]
\[
K_R(p;q_1,\ldots,q_5) = \int_X K(p;q_1,q_2,u)K(u;q_3,q_4,q_5)\,d\mu(u).
\]
With defect kernel $\Phi_K = K_L-K_R$,
\begin{align*}
A_K(f_1,\ldots,f_5)(p) &= T_K(T_K(f_1,f_2,f_3),f_4,f_5)(p) - T_K(f_1,f_2,T_K(f_3,f_4,f_5))(p)\\
&= \int_{X^5}\Phi_K(p;q_1,\ldots,q_5)\prod_{j=1}^5 f_j(q_j)\,d\mu^{\otimes 5}.
\end{align*}
Define the associator energy $\rho_A(K) := \lVert\Phi_K\rVert_{L^2(X^6)}^2$.

\begin{proposition}
If $\Phi_K=0$ a.e., the product is associative under the stated
convention; under suitable density/integrability conditions, if the
associator vanishes on all products of test functions then
$\Phi_K=0$ a.e.\ (converse also holds).
\end{proposition}

\section{Variational program}

Let $\theta$ parametrize a law $\mu_\theta$, kernel $K_\theta$, and
possibly a projector $P_\theta$. A general energy is
\[
\mathcal E(\theta) = \lambda_A\mathcal E_A + \lambda_C\mathcal E_C
+ \lambda_{\mathrm{cyc}}\mathcal E_{\mathrm{cyc}} + \lambda_{\mathrm{FI}}\mathcal E_{\mathrm{FI}}
+ \lambda_R\mathcal E_R + \lambda_P\mathcal E_P,
\]
e.g.\ $\mathcal E_A = \mathbb E\lVert A_{\mu_\theta}(X)\rVert^2$,
$\mathcal E_C = \mathbb E\lVert(I-P_\theta)\mu_\theta(P_\theta X_1,\ldots,P_\theta X_a)\rVert^2$,
$\mathcal E_P = \lVert P_\theta^2-P_\theta\rVert_F^2 + \lVert P_\theta-P_\theta^*\rVert_F^2$.
Parametrizing $P=QQ^*$, $Q^*Q=I_r$ places $Q$ on the Stiefel manifold
$\operatorname{St}(d,r)=\{Q\in\mathbb K^{d\times r}: Q^*Q=I_r\}$, with
tangent space $T_Q\operatorname{St}(d,r)=\{Z: Q^*Z+Z^*Q=0\}$; for a
Euclidean gradient $G$, the Riemannian gradient is
$\operatorname{grad}\mathcal E(Q) = G-Q\,\operatorname{sym}(Q^*G)$.

\begin{remark}
Finding a numerical minimum does not prove it is global, nor that the
recovered subspace is canonical.
\end{remark}

\section{Finite cohomology and truncated Hodge theory}
\label{sec:xxiii}

\subsection{Cochain complex}

Let $C^0\xrightarrow{d_0}C^1\xrightarrow{d_1}\cdots$ be a finite
complex with $d_{p+1}d_p=0$, and
$H^p(C,d) = \ker d_p/\operatorname{im}d_{p-1}$.

\subsection{Descent of an operator}

A graded operator $T$ induces $T_*:H^p\to H^p$ if
$T(\ker d_p)\subseteq\ker d_p$ and
$T(\operatorname{im}d_{p-1})\subseteq\operatorname{im}d_{p-1}$; a
sufficient condition is $Td=dT$, since then $x\sim x+dy \Rightarrow
Tx+d(Ty)$, so $[Tx]$ is well defined.

\subsection{Hodge compatibility}

With inner products and adjoint $d^*$, let
$\Delta = dd^*+d^*d$. If $[T,d]=0$ and $[T,d^*]=0$ then $[T,\Delta]=0$,
so $T$ preserves the harmonic subspaces $\mathcal H^p=\ker\Delta_p$.

\subsection{Spectral truncation}

On a compact manifold, let $\Pi_{\le\Lambda}$ be the spectral projector
of the Hodge Laplacian and $V_\Lambda=\ran\Pi_{\le\Lambda}$. For
curried operators $T^K_{e_i,e_j}$, define
\[
\mathcal E_d^\Lambda(K) = \sum_{i,j}\Bigl(
\lVert[d,T^K_{e_i,e_j}]\Pi_{\le\Lambda}\rVert_{\mathrm{HS}}^2
+\lVert[d^*,T^K_{e_i,e_j}]\Pi_{\le\Lambda}\rVert_{\mathrm{HS}}^2\Bigr).
\]
If $\mathcal E_d^\Lambda(K)=0$, the operators commute with $d,d^*$
within the truncation and descend to truncated cohomology. The
proposed total functional is
$\mathcal E_{\mathrm{total}}^\Lambda(K) = \mathcal E_{\mathrm{assoc}}^\Lambda(K)+\eta\,\mathcal E_d^\Lambda(K)$.

\begin{remark}
Being a local minimum does not by itself imply
$\mathcal E_d^\Lambda=0$; that conclusion needs an additional
control/surjectivity hypothesis on variations of the defect map.
\end{remark}

\section{Markov operator and induced Laplacian}

Suppose a declared kernel contraction yields a symmetric nonnegative
function $W_K(p,s)=W_K(s,p)\ge0$. Set
$d_K(p)=\int_X W_K(p,s)\,d\mu(s)$ and, where $d_K(p)>0$,
$P_K(p,s)=W_K(p,s)/d_K(p)$. With $d\nu(p)=d_K(p)\,d\mu(p)$,
$(P_Kf)(p)=\int_X P_K(p,s)f(s)\,d\mu(s)$, and $\Delta_K = I-P_K$,
$P_K$ is self-adjoint on $L^2(X,\nu)$ and
\[
\langle f,\Delta_K f\rangle_{L^2(\nu)}
= \frac12\iint \lvert f(p)-f(s)\rvert^2 W_K(p,s)\,d\mu(p)\,d\mu(s),
\]
so $\Delta_K\ge0$. If $e^{-t\Delta_K}$ is trace class, let
$\Theta(t)=\Tr(e^{-t\Delta_K})$. If $\Theta(t)\sim Ct^{-d_s/2}$ as
$t\downarrow0$, the spectral dimension is
$d_s = -2\lim_{t\downarrow0}\tfrac{d\log\Theta(t)}{d\log t}$.

\begin{remark}
Existence of this limit is an additional hypothesis, not an automatic
consequence of the ternary kernel.
\end{remark}

\section{Multiscale structure}
\label{sec:xxiv}

Let $(V_N,\mu_N,P_N)$ be a family indexed by resolution $N$, with
prolongation maps $J_N^M:V_N\to V_M$ and restriction maps
$R_M^N:V_M\to V_N$. Define the law transport defect
$\mathcal D_\mu^{N,M} = J_N^M\mu_N - \mu_M(J_N^M\cdot,\ldots,J_N^M\cdot)$
and the projector transport defect
$\mathcal D_P^{N,M} = P_MJ_N^M - J_N^MP_N$. Subspace distances may be
measured via $\lVert P_M - J_N^MP_NR_M^N\rVert$ or principal angles.

\begin{remark}
A finite decreasing sequence of such values does not by itself
establish a continuum limit. That would additionally require explicit
topologies, uniform bounds, compactness control, convergence of laws,
convergence of projectors, stability of the identities, and
identification of the limiting operator. The present computational
results are correctly limited to finite observations of transport, not
a continuum claim.
\end{remark}

\section{The complete canonical object}

The finite mathematical structure can be summarized as
\[
\mathfrak S = \Bigl(\mathcal T,\ \{V_\tau,W_\tau,Q_\tau,P_\tau\}_{\tau\in\mathcal T},\ \{\mu_v\}_v,\ \{K_v\}_v,\ \{A_v\}_v,\ \mathcal E,\ T,\ \mathcal C,\ \mathcal Z\Bigr),
\]
with $\mathcal T$ the types; $V_\tau,W_\tau$ ambient/reduced spaces;
$Q_\tau$ isometric inclusions; $P_\tau$ projectors; $\mu_v$ multilinear
laws; $K_v$ tensors or kernels; $A_v$ declared composition defects;
$\mathcal E$ functionals; $T$ the composition tree/forest; $\mathcal C$
error certificates; $\mathcal Z$ proof/exactness/evidence artifacts.
The canonical chain is
\begin{align*}
\text{typed space} &\to \text{multilinear law} \to \text{associator}
\to \text{closure leakage} \to \text{projector}\\
&\to \text{reduced law}
\to \text{recursive tree} \to \text{tangent-normal calculus}
\to \text{global certificate}.
\end{align*}

\section{Exact mathematical status}
\label{sec:xxvi}

\subsection{Proved in the finite core}
\[
(E_T^{\mathrm{amb}})^2 = (E_T^P)^2+(E_T^N)^2,
\qquad
E_T^{\mathrm{red}}=E_T^P,
\qquad
D_v = r_v+\sum_{\varnothing\neq S}\mu_v(y^S),
\]
\[
E_T^{\mathrm{amb}} \le k\rho M^{k-1}L_T,
\qquad
E_T^{P}=E_T^{\mathrm{red}} \le (k-1)\rho M^{k-1}L_T,
\]
\[
R_{\mathrm{std}}(x,y)z = A(y,x,z)-A(x,y,z),
\qquad
K\in L^2(X^{a+1}) \Longrightarrow T_K \text{ is bounded multilinear},
\]
\[
Td=dT \Longrightarrow T \text{ descends to cohomology}.
\]

\subsection{Proved, but of limited scope}
Soundness of the mixed-mask recurrence; the path-sum identity for the
declared recurrence; the optimal order within the declared telescoping
family; the triangle bound for signed forests; stability under tensor
perturbation; self-adjointness of the induced Laplacian under stated
symmetry and positivity.

\subsection{Open}
The exact value of $C_T^P(\eta)$; the exact constant for $k=2$;
per-topology constants for $k=3$; sharpness of $k-1$; sharpness of the
associator's coefficient $2$; cancellation-aware certificates;
informational optimality of the mixed-mask calculus; the continuum
limit; membership in $\Psi^0$; microlocal regularity; $D$-modules and
the Riemann--Hilbert correspondence; algebraization of the limiting
projector; theorem-to-theorem novelty.

\section{Final synthesis}

The defensible mathematical formalization of SEION is a finite theory
of typed multilinear laws, approximately invariant projectors, and
nodewise error propagation on composition trees. Its current main
theorem is
\[
E_T^{P} \le (k-1)\rho M^{k-1}L_T,
\]
and its central mechanism is that internal leakage can return to
tangent, but the root's own normal leakage is eliminated exactly
(Theorem~\ref{thm:orthogonality}). The kernel, cohomology, Hodge, and
spectral-operator extensions are mathematically formulated in a bounded
or truncated regime. The passage to the continuum, microlocal geometry,
and algebraization remain a conditioned program, not a closed theory.

\end{document}
